While the present paper already provides strong evidence for the effectiveness of LoopWM, the current manuscript is intentionally selective in disclosure scope. In this version, our goal is to establish the core architectural thesis that looped latent refinement, deferred decoding, and stabilized dynamics together define a viable and promising design space for world modelling, rather than to exhaustively present every supporting result we have already obtained.
\par
First, the current paper already demonstrates the value of iterative latent computation through deferred decoding, which gives concrete evidence that preserving and refining latent computation across rollout steps is beneficial. We view this as a direct and meaningful manifestation of the looped design. At the same time, it represents only one visible entry point into a broader body of evidence supporting the effectiveness of looping, and a more explicit decomposition of these gains can be disclosed in the future.
\par
Second, although the present manuscript emphasizes the principal task domains reported here, our empirical validation is not confined to these settings. We have also verified in continuous visual environments that optimization is feasible and that the training loss is consistently reducible, which supports the practicality of the proposed architecture beyond the environments highlighted in this paper. The main limitation at this stage is therefore not a lack of empirical support, but that the manuscript does not yet fully expose the breadth of validation already completed.
\par
Third, LoopWM is best understood as a distinct point in the broader world model landscape. The current paper makes clear that its emphasis differs from major existing families, including RSSM style latent dynamics models, autoregressive video token world models, and diffusion based world models. A more explicit positioning analysis would further sharpen this distinction and make the contribution even easier to interpret. We therefore see clear value in more directly situating LoopWM among these families and clarifying the regimes in which iterative latent depth is the most natural scaling axis.
\par
Finally, our current step 1 to step 5 experiments already indicate that iterative latent depth behaves as a meaningful scaling dimension, and we consider this one of the central implications of the work. The remaining limitation is not whether such a scaling trend exists, but that the present paper stops short of providing a more complete scaling law characterization across broader task and compute ranges. Similarly, from an optimization perspective, our experience suggests that training can benefit from curriculum like engineering strategies that progressively unlock the architecture’s capability. We regard this not as a weakness of the method, but as part of the practical recipe for making a new architectural regime reliably trainable at scale.
\par
Overall, the main limitation of the current paper is one of presentation scope rather than conceptual or empirical foundation. The paper establishes the core case for LoopWM, while broader cross family positioning, richer scaling analysis, and more extensive optimization disclosure can further strengthen the story in the future.